@@M0@@, which is plainly unique.

\sgaStableTarget{sga6.e05.prop.4-8}\textbf{Proposition 4.8.} \uline{Let} @@M1@@ \uline{be an integer} @@M2@@; \uline{in the free} @@M3@@-\uline{ring generated} \uline{over} @@M4@@ \uline{by an indeterminate} @@M5@@, \uline{the ideal generated by the} @@M6@@, @@M7@@, \uline{is} \uline{a} @@M8@@-\uline{ideal}.

Put @@M9@@, and let @@M10@@ be the ideal generated by the @@M11@@ for @@M12@@; put @@M13@@. We must show that the composite homomorphism: @@M14@@ vanishes on @@M15@@, i.e. that @@M16@@, @@M17@@, has image @@M18@@ in @@M19@@. Now @@M20@@, and @@M21@@ is a @@M22@@-homomorphism; it therefore suffices to show that for every ring @@M23@@ and every series @@M24@@ whose coefficients vanish in degrees @@M25@@, @@M26@@ when @@M27@@. But to prove this last property it suffices to prove it for @@M28@@, and @@M29@@. Embedding @@M30@@ in @@M31@@ by means of symmetric functions, we can write: @@M32@@, hence
@@M33@@
which has degree @@M34@@ in @@M35@@.

\sgaStableTarget{sga6.e05.cor.4-9}\textbf{Corollary 4.9.} \uline{The} @@M36@@-\uline{ring generated over} @@M37@@ \uline{by the generator} @@M38@@ \uline{subject to the relations} @@M39@@ \uline{for} @@M40@@, \uline{is isomorphic to} @@M41@@.

4.9.1. We would have the same statements on replacing the @@M42@@ by the @@M43@@. We can also replace @@M44@@ by a family of indeterminates @@M45@@, subject to the relations @@M46@@ for @@M47@@, where the @@M48@@ are integers @@M49@@; the resulting @@M50@@-ring is then @@M51@@.

\sgaStableTarget{sga6.e05.para.4-10}4.10. Let @@M52@@ be a binomial ring, and @@M53@@ an augmented @@M54@@-@@M55@@-algebra. We denote by @@M56@@ the free @@M57@@-ring generated by @@M58@@ over @@M59@@;
